% !TeX program = lualatex
\documentclass{linearspace-report}

\usepackage[backend=biber,style=numeric-comp,sorting=none]{biblatex}
\addbibresource{templates/shared-data/references.bib}

\LSSetup{
  title={Engineering Analysis and Verification Report},
  subtitle={A professional template for mathematical design documentation},
  short-title={Engineering Analysis Report},
  author={Engineering Team},
  organization={Linear Space},
  document-id={LS-ENG-0001},
  revision={A},
  status={Draft for Review},
  date={\today},
  project={Example Spacecraft Subsystem},
  client={Internal Template Demonstration}
}

% Safe default: this repository contains no controlled information. See
% docs/MARKINGS.md before enabling any controlled or classified profile.
\LSMarkingSetup{
  profile=uncontrolled,
  classification=none,
  controls={},
}

\begin{document}

\pagenumbering{roman}
\LSMakeReportTitle

\begin{LSRevisionHistory}
  A & 2026-08-30 & Engineering Team & Initial template release. \\
\end{LSRevisionHistory}

\chapter*{Executive Summary}
\addcontentsline{toc}{chapter}{Executive Summary}
This document demonstrates the intended structure for engineering analyses,
design descriptions, interface documents, verification reports, and manuals.
The numerical values and data are synthetic; replace them with traceable
project evidence.

\begin{LSKeyPoint}[Decision-ready conclusion]
State the engineering conclusion first, quantify its margin, identify the
limiting assumption, and name the evidence required to close any residual risk.
\end{LSKeyPoint}

\tableofcontents
\listoffigures
\listoftables
\clearpage
\pagenumbering{arabic}

\chapter{Purpose and Scope}

This report defines the analysis boundary, mathematical model, verification
method, and acceptance evidence for an illustrative attitude-estimation
subsystem. Classical stochastic-estimation references such as
\textcite{maybeck1979} provide the surrounding theory; project-specific
assumptions must still be justified and controlled.

\section{Objectives}

\begin{LSRequirements}
  \item\label{req:accuracy} The steady-state three-axis estimation error shall
    remain below \qty{1.5}{\milli\radian} at the defined confidence level.
  \item\label{req:rate} The implementation shall produce an estimate at
    \qty{10}{\hertz} or faster under the target processor budget.
  \item\label{req:traceability} Every reported result shall identify its model,
    input data, software revision, and acceptance criterion.
\end{LSRequirements}

\section{Document Conventions}

Vectors are bold lowercase, matrices are bold uppercase, and coordinate frames
appear as right superscripts. SI units are typeset with \texttt{siunitx}; figures,
tables, equations, and requirements are cross-referenced with \texttt{cleveref}.

\chapter{Mathematical Model}

\section{State Propagation}

Let the attitude-error state be
\begin{equation}
  \vect{x}
  =
  \begin{bmatrix}
    \delta\vect{\theta}^{\transpose} & \vect{b}^{\transpose}
  \end{bmatrix}^{\transpose}
  \in \Rset^{6},
  \label{eq:state}
\end{equation}
where \(\delta\vect{\theta}\) is the small-angle error and \(\vect{b}\) is
the gyro-bias estimate. A discrete linearization has the form
\begin{align}
  \vect{x}_{k+1} &= \mat{\Phi}_k\vect{x}_k + \mat{G}_k\vect{w}_k, \\
  \mat{P}_{k+1}^{-} &=
    \mat{\Phi}_k\mat{P}_{k}^{+}\mat{\Phi}_k^{\transpose}
    + \mat{G}_k\mat{Q}_k\mat{G}_k^{\transpose}.
  \label{eq:covariance}
\end{align}

The covariance update in \cref{eq:covariance} is only defensible when the
process-noise model, units, sampling interval, and frame conventions are all
explicit. For a full spacecraft-dynamics treatment, see \textcite{schaub2018}.

\begin{LSWarning}[Model limitation]
The illustrative equations omit normalization, finite-precision effects, and
sensor timing error. Do not convert a teaching model into a design claim without
closing those omissions.
\end{LSWarning}

\section{Geometry and Frames}

\begin{figure}[htbp]
  \centering
  \begin{tikzpicture}[
    node distance=1.0in,
    box/.style={draw=LSRule,fill=LSIce,rounded corners=2pt,minimum width=1.55in,minimum height=0.55in,align=center,font=\sffamily\small},
    flow/.style={-{Latex[length=2.5mm]},very thick,draw=LSBlue}
  ]
    \node[box] (sensor) {Sensor samples};
    \node[box,right=of sensor] (prop) {State propagation};
    \node[box,right=of prop] (update) {Measurement update};
    \node[box,below=0.55in of prop] (evidence) {Residuals and evidence};
    \draw[flow] (sensor) -- (prop);
    \draw[flow] (prop) -- (update);
    \draw[flow] (update) |- (evidence);
    \draw[flow] (evidence.north) -- (prop.south);
  \end{tikzpicture}
  \caption{Traceable estimation and verification flow.}
  \label{fig:flow}
\end{figure}

\chapter{Verification Evidence}

\section{Acceptance Matrix}

\begin{table}[htbp]
  \centering
  \caption{Illustrative verification matrix.}
  \label{tab:verification}
  \begin{tabularx}{\linewidth}{@{}p{0.95in}Y p{1.00in}p{0.75in}@{}}
    \toprule
    \textbf{Requirement} & \textbf{Method and evidence} & \textbf{Criterion} & \textbf{Status} \\
    \midrule
    \cref{req:accuracy} & Monte Carlo simulation with archived seeds, input set,
      and percentile summary & \(<\qty{1.5}{\milli\radian}\) & Example \\
    \cref{req:rate} & Hardware-in-the-loop timing profile with worst-case load &
      \(\geq\qty{10}{\hertz}\) & Pending \\
    \cref{req:traceability} & Configuration audit of report, data, code, and build
      identifiers & Complete record & Pending \\
    \bottomrule
  \end{tabularx}
\end{table}

\section{Plotting Quantitative Results}

\begin{figure}[htbp]
  \centering
  \begin{tikzpicture}
    \begin{axis}[
      linearspace,
      width=0.88\linewidth,
      height=2.75in,
      xlabel={Time (s)},
      ylabel={Angular error (mrad)},
      xmin=0,xmax=10,
      ymin=0,ymax=2.0,
      legend columns=3,
      legend to name=verificationlegend,
    ]
      \addplot table[x=time,y=error,col sep=comma]
        {templates/shared-data/estimation-error.csv};
      \addlegendentry{Observed error}
      \addplot table[x=time,y=one_sigma,col sep=comma]
        {templates/shared-data/estimation-error.csv};
      \addlegendentry{One-sigma envelope}
      \addplot table[x=time,y=requirement,col sep=comma]
        {templates/shared-data/estimation-error.csv};
      \addlegendentry{Requirement}
    \end{axis}
  \end{tikzpicture}
  \par\vspace{0.08in}\ref{verificationlegend}
  \caption{Synthetic convergence data rendered directly from CSV.}
  \label{fig:convergence}
\end{figure}

\Cref{fig:convergence} demonstrates the preferred workflow: keep source data
machine-readable, render vector plots in LaTeX, label units, and state the
meaning of every curve. Do not infer compliance from appearance alone.

\chapter{Conclusions and Release Checklist}

\begin{LSKeyPoint}[Recommended disposition]
Replace this paragraph with the decision, quantitative margin, open risks,
owners, and due dates. Separate verified facts from engineering judgment.
\end{LSKeyPoint}

Before release, confirm that:
\begin{itemize}
  \item all requirements and interfaces are traceable;
  \item plots are regenerated from controlled data, not screenshots;
  \item units, coordinate frames, and uncertainty conventions are explicit;
  \item classification, CUI, export-control, proprietary, and distribution
    markings were assigned by the proper authority; and
  \item the source, build directory, logs, auxiliary files, and PDF are stored
    only in an authorized environment.
\end{itemize}

\printbibliography[heading=bibintoc]

\appendix
\chapter{Reusable Portion-Marking Examples}

The following commands are intentionally shown as code rather than emitting
controlled markings in this public template:
\begin{lstlisting}[language={[LaTeX]TeX}]
\LSU{Unclassified paragraph text.}
\LSCUI[SP-CTI]{Controlled technical information.}
\LSS{Secret portion text.}
\LSTS{Top Secret portion text.}
\end{lstlisting}
See \texttt{docs/MARKINGS.md} before using any of them.

\end{document}
